% =============================================================================
%  Sección 5: Conclusiones
% =============================================================================

\section{Conclusiones}

A lo largo de este ensayo, se ha analizado la definición y el alcance de los humedales en el ámbito metropolitano, 
lo cual se traduce en tres ejes estratégicos: social, arquitectónico y biológico. Estos tres pilares constituyen 
la base de la sostenibilidad urbana a largo plazo. Los ecosistemas hidrofílicos son de suma importancia para la 
captación de agua, el refugio natural de especies y la concentración de actividades económicas para los pobladores locales, 
lo que los convierte en entes vivos de importancia estratégica para la resiliencia de la Ciudad de México \cite{paot2006xochimilco, sedema2022programa}. 

Las conclusiones de este estudio se presentan bajo los siguientes ejes rectores:

\subsection{Eje Social: Identidad, Resistencia y la Reivindicación del Poblador}

Los humedales representan un nodo sociocultural donde la identidad se debate entre la preservación y la exclusión. 
El sistema de chinampas en Xochimilco es el ejemplo más claro de una interacción equilibrada entre la especie humana y el ambiente, 
favoreciendo el crecimiento de especies endémicas \cite{paot2006xochimilco}. No obstante, la historia urbana reciente muestra un 
quiebre en este equilibrio; la política de extracción hídrica resultó en la degradación de manantiales y su sustitución por aguas 
residuales \cite{paot2006xochimilco}. 

Como consecuencia, poblados como Caltongo y San Gregorio Atlapulco han sido estigmatizados como focos de infección, ignorando 
que la degradación no deriva de los pobladores, sino de políticas públicas fallidas y desinterés institucional \cite{rojas2020deterioro}. 
Esta desatención en servicios básicos obliga a las comunidades a adoptar posturas herméticas y desplazarse hacia zonas poco aptas para el 
hábitat humano, como se observa en la zona lacustre de la avenida Nuevo León \cite{rojas2020deterioro}.

\subsection{Eje de Planeación: El Fracaso de la Expansión Horizontal y el Acceso a la Ciudad}

El crecimiento de la Zona Metropolitana del Valle de México hacia sus periferias revela una incapacidad de planeación vertical, 
forzando a las zonas limítrofes a convertirse en áreas de exclusión. Mientras que en Xochimilco la infraestructura gris como el puente de 
Cuemanco amenaza con fragmentar el territorio \cite{segob2020proposicion}, en Aragón la urbanización temprana (1970-1990) consumió gran parte de su identidad original. 

Sin embargo, el caso de Aragón demuestra que es posible recuperar estos ecosistemas mediante humedales artificiales que funcionan como atractivos turísticos 
y nodos de saneamiento \cite{unam2020aragon, sedema2022programa}. El reto actual reside en que los pobladores de la periferia dependen de los Distritos Centrales 
de Negocios (CBD), como el eje Reforma, para acceder a servicios de salud, trabajo y burocracia que no se encuentran en sus alcaldías \cite{rojas2020deterioro}.

\subsection{Eje de Movilidad: El Transporte como Ente Regulador y no Destructor}

La palabra "progreso" se ha vinculado históricamente con la destrucción de ecosistemas, pero bajo una visión de sustentabilidad, 
debe transitar hacia la regulación. El transporte masivo es una necesidad básica cuya ineficiencia afecta la escala macro de la metrópoli. 
Bajo el \textit{Modelo VFT}, la propuesta de un corredor circular sobre el Anillo Periférico debe actuar como una frontera estructural \cite{garibaldi2023vft}. 

Al gestionar matemáticamente la fuerza capilar del transporte y consolidar los flujos en nodos optimizados, la arquitectura vial puede fungir como una barrera 
artificial que detenga la dispersión de transeúntes hacia áreas naturales protegidas \cite{garibaldi2023vft}. Entender que un transporte eficiente define los 
límites entre lo rural, lo urbano y lo ambiental es clave para un progreso sistemático que concilie los intereses del gobierno, la ciudadanía y la empresa, 
protegiendo así el patrimonio biológico de la ciudad \cite{garibaldi2023vft, sedema2022programa}.
